% GENERATED FILE. Edit canonical lessons, then run npm run generate:books.
% canonical-source-hash: fnv1a64:a4b42598f9956103
% canonical-lessons: KA-C27-sanje

\chapter{Evening}
\label{ch:ka-evening}

\begin{chapteropening}
\textbf{By the end of this chapter:} \emph{I can say \textquotedblleft{}evening\textquotedblright{} in Kannada.}

Say \kn{ಸಂಜೆ}, follow it back through Maharashtri Prakrit to Sanskrit saṃdhyā, and set it beside Hindi's \dv{साँझ}, the same root down a different Prakrit road.
\end{chapteropening}

\section[\kn{ಸಂಜೆ}]{\kn{ಸಂಜೆ} (sañje) — a Sanskrit word that arrived already worn down}

\label{lesson:KA-C27-sanje}



\begin{quote}
\kn{ದಿನ}, \kn{ರಾತ್ರಿ}, and \kn{ಮಧ್ಯಾಹ್ನ} all arrived in Kannada whole, unworn, straight from Sanskrit. This word took a very different route — and it lands somewhere genuinely surprising.
\end{quote}

\begin{cousinweb}
\textbf{\kn{ಸಂಜೆ}} (\textbf{sañje}) — \textquotedblleft{}\textbf{evening}\textquotedblright{} — traces to \textbf{Sanskrit} \textbf{\kn{ಸಂಧ್ಯಾ}} (\emph{saṃdhyā}, \textquotedblleft{}\textbf{twilight},\textquotedblright{} literally \textquotedblleft{}the junction of day and night\textquotedblright{}) — but it didn't arrive in Kannada as a fresh, whole tatsama the way \emph{dina}, \emph{ratri}, and \emph{madhyāhna} did. Instead, it came by way of \textbf{Maharashtri Prakrit} \textbf{\kn{ಸಂಝಾ}} (\emph{sañjhā}) — Sanskrit's word, already worn down by centuries of everyday Middle Indo-Aryan speech, \emph{then} borrowed into Kannada in that already-eroded shape.
\end{cousinweb}

\begin{cousinweb}
Here's the genuinely surprising part: Hindi's own evening word, \textbf{\dv{साँझ}} (\emph{sāñjh}), comes from the \textbf{exact same} Sanskrit \textbf{\kn{ಸಂಧ್ಯಾ}} — but by a \textbf{different} route. Hindi \emph{inherited} it, generation to generation, through \textbf{Śauraseni Prakrit} (a different, sister Middle Indo-Aryan dialect than the Maharashtri Prakrit Kannada borrowed from). Two separate erosion paths — one an unbroken inheritance within the Indo-Aryan family, one a loan crossing into a completely different language family — and they land on nearly \textbf{identical-sounding} modern words, \emph{sañje} and \emph{sāñjh}, purely because they share the same Sanskrit starting point.
\end{cousinweb}

\subsection*{Your turn}
Say these aloud:

\begin{itemize}
\raggedright
  \item \textquotedblleft{}sañje\textquotedblright{} — \textquotedblleft{}evening,\textquotedblright{} from Sanskrit sandhyā via Maharashtri Prakrit
  \item the contrast with dina/ratri/madhyāhna — those arrived whole; sañje arrived already worn down
  \item the striking convergence — Hindi's \dv{साँझ}, same Sanskrit root, different Prakrit path, nearly the same sound today
\end{itemize}

\begin{tcolorbox}[breakable,colback=teal!4,colframe=teal!35!black,title={Before you move on}]
Does \kn{ಸಂಜೆ} arrive in Kannada as a fresh, whole Sanskrit tatsama, like \emph{dina} and \emph{ratri} did? (\textbf{No} — it arrived already worn down, via Maharashtri Prakrit \kn{ಸಂಝಾ}.) What Hindi word shares \kn{ಸಂಜೆ}'s exact Sanskrit root, and by what different path did it get there? (\textbf{\dv{साँझ}} (sāñjh) — inherited via a different sister Prakrit, Śauraseni, rather than borrowed via Maharashtri.)
\end{tcolorbox}
